% -*- Mode: Latex -*-

%  Zusammenfassung

% Zu einer runden Arbeit gehört auch eine Zusammenfassung, die
% eigenständig einen kurzen Abriß der Arbeit gibt. Eine halbe bis ganze
% DINA4 Seite ist angemessen. Dafür läßt sich keine Gebrauchsanweisung
% geben (für irgendetwas müssen die Betreuer ja auch noch da
% sein).

Modern hardware increasingly combines processor cores of varying types and speeds on a single chip (asymmetric multiprocessing, or AMP).
However, the L4Re operating system framework does not provide CPU load balancing.
Threads need to be statically assigned to cores by the applications themselves.
This approach aligns with the microkernel philosophy of separating mechanism from policy, but it leaves the question of AMP-aware, dynamic thread placement unaddressed.

This thesis proposes a user-level load balancing service for L4Re on x86-64 that distributes threads across physical cores during runtime.
It respects application-supplied constraints, such as thread pinning, core-sharing restrictions, and the placement requirements of virtual machine vCPUs.
The design adheres to L4Re's proxy pattern: a scheduling proxy intercepts scheduler calls, per-CPU monitoring threads track system state, and a commander thread performs periodic rebalancing based on virtual load (vload), a preemption-based metric capturing both CPU utilization and saturation.
AMP awareness is achieved by scaling the vload of each CPU according to its measured compute performance.
Realizing the design required extending the L4Re microkernel with three new interfaces for capability comparison, thread probing, and pinned scheduling.

The resulting load balancer is evaluated against the Linux EEVDF scheduler using ported CPU-bound benchmarks across baseline, scalability, and parallel-load scenarios while its own monitoring overhead is measured directly.
It is imperative to implement periodic rebalancing procedures, as a non-rebalancing variant trails by a factor of two to four due to a pile-on effect from stale load information.
The rebalancing configurations remain within a narrow range of each other, irrespective of the monitoring interval.
They introduce a monitoring overhead of less than 1\,\% and approach EEVDF's throughput within approximately 2--7\,\% percent.
This demonstrates that AMP-aware load balancing is achievable as a policy-separated userland service at a modest cost relative to a mature monolithic scheduler.

%%% Local Variables:
%%% TeX-master: "diplom"
%%% End:


